%--------------------------------------------------------%
% ARTICLE STRUCTURE GUIDELINES
%--------------------------------------------------------%
%
% Please structure your manuscript according to the category of your submission.
% Below is a general outline for different types of articles.
%
% For "Short Narrative Reviews":
% - Introduction
% - Methodology
% - Thematic Overview
% - Discussion
% - Conclusion
%
% For "Applied AI Exploration Papers":
% - Introduction
% - Project Description
% - Implementation and Results
% - Discussion and Potential Impact
% - Conclusion
%
% For "Original Research Articles":
% - Must start with an Introduction and end with a Conclusion. The structure in between these sections is flexible.
%
% For detailed instructions on specific sections and the structure of the papers,
% please refer to the "Guideline for Authors" on our website:
% https://maikron.org/jaica/index.php/ojs/guideline-for-authors
%
% Please ensure each section is clearly marked with \section{Section Title}.
%
%--------------------------------------------------------%
%  HOW TO CITE, CROSS-REFERENCE, AND ADD FIGURES/TABLES
%--------------------------------------------------------%
%
% Citations: add an entry to 'user/biblio.bib', then cite it with
%   \cite{yourkey}
%
% Figures: place image files in 'user/figures/', then include them with
%   \begin{figure}[htbp]
%   \centering
%   \includegraphics[width=1\textwidth]{user/figures/yourimage.png}
%   \caption{Your caption here.}
%   \label{fig:yourlabel}
%   \end{figure}
%   ... and refer to it in text with Figure~\ref{fig:yourlabel}
%
% Tables: use the tabularx environment already loaded in the preamble, e.g.
%   \begin{table}[H]
%   \centering
%   \caption{Your caption here.}
%   \label{tab:yourlabel}
%   \begin{tabularx}{0.95\textwidth}{XXc}
%   \hline
%   \textbf{Column A} & \textbf{Column B} & \textbf{Column C} \\
%   \hline
%   Row 1 value & Row 1 value & Row 1 value \\
%   \hline
%   \end{tabularx}
%   \end{table}
%   ... and refer to it in text with Table~\ref{tab:yourlabel}
%
%--------------------------------------------------------%
%	BODY TEXT
%--------------------------------------------------------%

\section{Introduction}
% Present the topic, its importance, the background needed to understand
% it, the scope of your paper, and your objectives or research questions.

[Write your introduction here.]

\section{Methodology}
\label{sec:methodology}
% Describe your approach, data sources, search strategy, inclusion/exclusion
% criteria, or experimental design, as appropriate for your type of paper.

[Write your methodology here.]

\section{Other needed sections}
\label{sec:thematic-overview}
% For reviews: synthesize what the literature shows, organized by theme.
% For research/applied papers, rename this section (e.g. "Results",
% "Implementation and Results") to match your paper type.

[Write the main body of your paper here.]

\section{Discussion}
\label{sec:discussion}
% Analyze and interpret your findings, situate them in the wider
% literature, and discuss limitations.

[Write your discussion here.]

\section{Conclusion}
\label{sec:Conclusion}
% Summarize your main findings, restate their significance, and suggest
% directions for future work.

[Write your conclusion here.]

%--------------------------------------------------------%
%  REQUIRED CLOSING STATEMENTS
%--------------------------------------------------------%
% The four statements below are typically required by the journal.
% Fill in each one; do not delete the sections themselves.

\section*{Ethics Statement}
[State whether the study involved human participants or animals and,
if so, how ethical approval was obtained. If not applicable, say so.]

\section*{CRediT authorship contribution statement}
% List each author and their contribution using CRediT roles
% (Conceptualization, Methodology, Software, Validation, Formal analysis,
% Investigation, Resources, Data curation, Writing - original draft,
% Writing - review & editing, Visualization, Supervision, Project
% administration, Funding acquisition). Example:
% \textbf{First Author Name:} Conceptualization, Methodology, Writing - original draft.
% \textbf{Second Author Name:} Supervision, Writing - review \& editing.

[List each author's contributions here.]

\section*{Declaration of Generative AI and AI-assisted technologies in the writing process}
% State whether, and how, generative AI tools were used in preparing the
% manuscript. If none were used, say so explicitly.

[State your use of generative AI tools here, or note that none were used.]

\section*{Declaration of competing interest}
% State any financial or personal relationships that could be seen as
% influencing the work, or declare that there are none.

[Declare competing interests here, or state that there are none.]
S